% ChkTeX's generic prose rules misclassify intentional TeX conventions used
% throughout this document: en-dashes, labelled equations, powers of grouped
% fractions, and punctuation in bibliographic titles.
% chktex-file 1
% chktex-file 2
% chktex-file 3
% chktex-file 8
% chktex-file 13
% chktex-file 24
% chktex-file 38
\documentclass[11pt]{article}

\usepackage[T1]{fontenc}
\usepackage[utf8]{inputenc}
\usepackage{lmodern}
\usepackage{microtype}
\usepackage[a4paper,margin=25mm,headheight=14pt]{geometry}
\usepackage{amsmath,amssymb,mathtools}
\usepackage{booktabs,longtable,tabularx,array,pdflscape}
\usepackage{xcolor}
\usepackage{pid-rs-report-tables}
\usepackage{enumitem}
\usepackage{xurl}
\usepackage{hyperref}

\ifdefined\pdfinfoomitdate%
  \pdfinfoomitdate=1
\fi
\ifdefined\pdftrailerid%
  \pdftrailerid{}
\fi
\ifdefined\pdfsuppressptexinfo%
  \pdfsuppressptexinfo=15
\fi

\hypersetup{%
  colorlinks=true,
  linkcolor=blue!55!black,
  citecolor=blue!55!black,
  urlcolor=blue!55!black,
  pdftitle={Foundational validity audit of shared-exclusions PID},
  pdfauthor={pid-rs project}
}

\PidConfigureSections
\PidSetRunningHeads{pid-rs foundational SxPID audit}{validity and compatibility}
\PidConfigureAbstract
\PidConfigureLinksAndListings
\setlength{\emergencystretch}{3em}
\setlist[itemize]{leftmargin=1.6em,itemsep=0.2em,topsep=0.4em}
\setlist[enumerate]{leftmargin=1.8em,itemsep=0.2em,topsep=0.4em}
\renewcommand{\arraystretch}{1.16}
\newcolumntype{Y}{>{\raggedright\arraybackslash}X}
\newcolumntype{P}[1]{>{\raggedright\arraybackslash}p{#1}}
\newcommand{\Sx}{\mathrm{sx}}
\newcommand{\Syn}{\operatorname{Syn}}
\newcommand{\evidence}[1]{\textbf{[#1]}}

\title{\PidReportTitle
  {Foundational Validity Audit of Shared-Exclusions PID}
  {Exact semantics, compatibility firewall, adversarial counterexamples,\\
   and a bounded executable audit}}
\author{pid-rs project}
\date{25 July 2026}

\begin{document}
\PidMakeTitle

\begin{abstract}
This report audits the categorical shared-exclusions partial information
decomposition (SxPID) of Makkeh, Gutknecht, and Wibral from its primitive event
semantics.  We rederive its lattice order, informative--misinformative split,
M\"obius inversion, and mutual-information reconstruction; distinguish proved
properties from constitutive semantic choices; and give exact finite
counterexamples to stronger interpretations.  We also reproduce the 2026
Lyu--Clark--Raviv relation witness.  A standard-library exact-rational checker
and a separate Rust regression show that their stipulated coordinate
recoverability descriptors are not SxPID atoms.  A kernel-checked Lean theorem
then isolates the missing transfer premise: an atom assignment must factor
through those descriptors before their collision can imply a reconstruction
impossibility for that assignment.  No fatal algebraic
contradiction is found in the native shared-exclusions construction.  The
valid conclusion is narrower: SxPID is a coherent signed, local, event-logical
decomposition under an explicit semantic and statistical contract, not a
universal ontology of informational parts.  Natural logarithms and nats are
used unless bits are stated explicitly.
\end{abstract}

\tableofcontents

\section{Executive decision}

The categorical shared-exclusions construction is mathematically valid in the
following precise sense.
\begin{enumerate}
  \item Its keyed source event is well-defined for every supported
        realization.
  \item Its local cumulative is an ordinary likelihood-ratio information
        value for that event.
  \item The antichain order, event containments, and M\"obius inversion give
        an exact signed decomposition.
  \item The greatest-node cumulative is local mutual information; averaging
        the atoms reconstructs joint source--target mutual information.
  \item Informative and misinformative component atoms are nonnegative.  Net
        atoms need not be.
  \item Source permutation symmetry, coordinatewise bijective re-encoding
        invariance, self-redundancy, and the target chain rule follow from the
        definition.
  \item The current Rust implementation reproduces the exact finite
        witnesses in this report, including the adversarial three-source
        relation witness in Section~\ref{sec:lcr}.
\end{enumerate}

No fatal algebraic contradiction was found in the Makkeh--Gutknecht--Wibral
construction or in the core theorems audited here.  This does \emph{not} show
that shared exclusions is the unique correct PID\@.  The following stronger
readings are false or unsupported:
\begin{itemize}
  \item the disjunction event is uniquely forced by Shannon theory;
  \item an atom is necessarily a nonnegative set-sized part, a common
        deterministic signal, a common target feature, a causal contribution,
        or a dependence strength;
  \item independent sources must have zero Sx redundancy;
  \item signed net atoms obey a general data-processing inequality;
  \item pairwise source--target marginals identify the decomposition;
  \item source and target roles are interchangeable;
  \item the changing keyed event is one fixed auxiliary random variable; or
  \item an observational PID atom establishes mechanism, responsibility,
        safety, deception, or decision authority.
\end{itemize}

\begin{quote}
\textbf{Shared exclusions is a coherent, signed, local, event-logical PID under
an explicit semantic contract.  It is not a universal ontology of
informational parts.}
\end{quote}

\section{Evidence discipline and audit protocol}

Every material result is classified as follows.
\begin{center}
\begin{tabularx}{\textwidth}{@{}lY@{}}
\toprule
\PidTableHeaderRow
Label & Meaning \\
\midrule
\evidence{P} & Published theorem or definition in a primary source. \\
\evidence{R} & Independent exact rederivation in this audit. \\
\evidence{X} & Exact finite counterexample or exact finite witness. \\
\evidence{B} & Bounded exhaustive or seeded computation; evidence, not a
universal proof. \\
\evidence{E} & Observation from the current \texttt{pid-rs} executable path. \\
\evidence{O} & Open question or unproved extension. \\
\bottomrule
\end{tabularx}
\end{center}

The protocol was to freeze the native primitive; derive the lattice and
reconstruction from events; separate semantics from numerical theorems; seek
minimal exact counterexamples; enumerate bounded empirical laws with rational
arithmetic; run independent executable witnesses; and apply the compatibility
firewall in Section~\ref{sec:firewall}.  Each conclusion was reviewed through
five lenses: definitional compatibility, theorem scope, executable refinement,
numerical/statistical implications, and adversarial counterexamples.

\subsection{Five-lens result ledger}

\begin{landscape}
\begingroup
\setlength{\tabcolsep}{3pt}
\scriptsize
\begin{longtable}{@{}P{0.10\linewidth}P{0.15\linewidth}P{0.17\linewidth}P{0.16\linewidth}P{0.17\linewidth}P{0.17\linewidth}@{}}
\toprule
\PidTableHeaderRow
Result & Definition & Theorem and assumptions & Executable refinement &
Numerical/statistical implication & Adversarial finding \\
\midrule
\endfirsthead
\toprule
\PidTableHeaderRow
Result & Definition & Theorem and assumptions & Executable refinement &
Numerical/statistical implication & Adversarial finding \\
\midrule
\endhead
F1 Keyed event & Native & Valid for supported finite keys & Rust scans exact
equality events & Rare events produce large logs & Unsupported keys have no
local occurrence \\
F2 Lattice/inversion & Native & Finite-poset inversion is exact & Two--four
source lattices implemented & Cancellation amplifies rounding & Labels do not
prove reconstruction \\
F3 Component signs & Native & Applies separately to $+$ and $-$, not net &
Exact bounded searches agree & Interval signs needed near zero & Exact negative
net atoms exist \\
F4 Averaging & Native & Linearity and top-node self-redundancy & Exact witnesses
reconstruct MI & Average is not MI of one fixed keyed event & Relation witness
still reconstructs \\
F5 REI/target chain & Native & Coordinatewise bijections; conditional chain law
& Relabeling/lattice tests & Identities hold only within numerical error &
Mixing/coarsening are outside REI \\
F6 Identity/LP & Comparison properties & Sx rejects both generally & COPY and
target-copy witnesses & Signed inference is mandatory & Independent sources
can have positive redundancy \\
F7 Disjunction & Native postulate & Logical consequence is valid; interpretation
is chosen & Code evaluates that choice & Another semantics is another estimand
& No uniqueness theorem follows \\
F8 Common signal & Stronger incompatible semantics & Not implied by Sx & No
common-variable object is computed & Do not call Sx G\'acs--K\"orner information
& Independent COPY is a counterexample \\
F9 Role/processing & Native values; stronger DPI incompatible & Direction is
built in & Witnesses executable & Preprocessing changes the estimand & Exact
role-swap/coarsening failures \\
F10 Identification & Full joint law required & Population functional is
identified by full finite law & Plug-in uses empirical joint PMF & Sparse
high-cardinality estimation is hard & Same pairwise marginals, different Sx \\
F11 Lyu--Clark--Raviv & Their atoms are not Sx atoms & Descriptor theorem is
internal; universal wording overbroad & Regression rejects substitution &
Relations change Sx atoms & Eight of eighteen atoms differ exactly \\
F12 Interpretation & Advisory statistic only & No causal or authority theorem &
Reports require estimand identity & UQ/selection are separate & Observationally
equivalent causal systems defeat attribution \\
\bottomrule
\end{longtable}
\normalsize
\endgroup
\end{landscape}

\section{Compatibility firewall}
\label{sec:firewall}

No external theorem is transferred to shared exclusions until its objects and
assumptions match Sections~\ref{sec:primitive}--\ref{sec:mobius}.

\small
\begin{longtable}{@{}P{0.19\textwidth}P{0.17\textwidth}P{0.23\textwidth}P{0.29\textwidth}@{}}
\toprule
\PidTableHeaderRow
Source or concept & Compatibility & Permitted use & Forbidden transplant \\
\midrule
\endfirsthead
\toprule
\PidTableHeaderRow
Source or concept & Compatibility & Permitted use & Forbidden transplant \\
\midrule
\endhead
Makkeh--Gutknecht--Wibral (2021) & Native & Definition, cumulatives,
component split, target chain, differentiability & Extending finite-discrete or
interior-domain theorems beyond their premises \\
Gutknecht--Wibral--Makkeh (2021) & Native conceptual foundation & Parthood and
formal-logic indexing & Calling constitutive principles uniquely forced
Shannon theorems \\
Matthias--Makkeh--Wibral--Gutknecht (2025) & Compatible abstract theorem &
LP--TCR--REI and LP--ID--REI incompatibilities & Calling either a defect unique
to SxPID or an algebraic inconsistency of Sx \\
Lyu--Clark--Raviv (2026), Definition 6/Lemma 4 & Incompatible atom semantics &
Coordinate-recoverability descriptors and relation loss & Calling
$H(U_\alpha)$ an SxPID atom or transferring their impossibility theorem without
the missing descriptor premise \\
Williams--Beer identity/local positivity & Comparison only & Report what Sx
accepts or rejects & Redefining Sx after a violation \\
Dependency coloring & Estimator layer only & Concentration under proved mutual
within-color independence & Calling it a ``colored PID'' measure or attributing
it to Wibral \\
\texttt{pid-rs} categorical kernel & Executable approximation & Regression and
refinement evidence & Treating tests as proof of semantics or exact-real
floating point \\
Project support-change theorem & Native averaged functional, project theorem &
Closed-simplex continuity on a fixed alphabet & Attributing it to the 2021 paper
or applying it to disappearing pointwise keys \\
\bottomrule
\end{longtable}
\normalsize

\section{Primitive reconstructed from first principles}
\label{sec:primitive}

Let
\[
Z=(S_1,\ldots,S_n,T)
\]
have law $p$ on a finite Cartesian-product alphabet.  Fix a supported key
\[
z=(s_1,\ldots,s_n,t),\qquad p(z)>0.
\]
For each nonempty collection $a\subseteq[n]$, define
\[
E_a(s)=\bigcap_{i\in a}\{S_i=s_i\}.
\]
For an antichain $\alpha$ of nonempty collections, define the inclusive
disjunction
\begin{equation}
A_\alpha(s)=\bigcup_{a\in\alpha}E_a(s).
\label{eq:event}
\end{equation}
Every branch contains the keyed source realization.  Therefore
$p(A_\alpha)>0$ and $p(T=t,A_\alpha)>0$.  The pointwise cumulative is
\begin{equation}
c_\alpha(z;p)
=\log\frac{p(T=t,A_\alpha)}{p(T=t)p(A_\alpha)}
=\log\frac{p(T=t\mid A_\alpha)}{p(T=t)}.
\label{eq:cumulative}
\end{equation}
It is a local likelihood ratio for the \emph{keyed} event, not yet an atom.
The informative and misinformative components are
\begin{equation}
c_\alpha^+(z;p)=-\log p(A_\alpha),
\qquad
c_\alpha^-(z;p)=\log\frac{p(T=t)}{p(T=t,A_\alpha)},
\label{eq:components}
\end{equation}
and $c_\alpha=c_\alpha^+-c_\alpha^-$.  Both components are nonnegative;
their difference can have either sign.

\paragraph{Finding F1 \evidence{P,R}.}
Equations~\eqref{eq:event}--\eqref{eq:components} are defined for every
supported finite key.  They assign no pointwise value to a zero-probability
realization.

\subsection{Antichain order}

Define
\begin{equation}
\alpha\preceq\beta
\quad\Longleftrightarrow\quad
\forall b\in\beta\;\exists a\in\alpha:\;a\subseteq b.
\label{eq:order}
\end{equation}
Since $a\subseteq b$ implies $E_b(s)\subseteq E_a(s)$,
\begin{equation}
\alpha\preceq\beta\quad\Longrightarrow\quad
A_\beta(s)\subseteq A_\alpha(s).
\label{eq:containment}
\end{equation}
Both $c^+$ and $c^-$ increase along the lattice.  Their difference need not.

\subsection{M\"obius atoms and exact reconstruction}
\label{sec:mobius}

Define atoms through the finite-poset zeta relation
\begin{equation}
c_\alpha^u(z;p)=\sum_{\beta\preceq\alpha}\pi_\beta^u(z;p),
\qquad u\in\{+,-,\mathrm{net}\},
\label{eq:zeta}
\end{equation}
where $\pi^{\mathrm{net}}=\pi^+-\pi^-$.  M\"obius inversion makes
Equation~\eqref{eq:zeta} unique after the cumulatives are fixed.  The greatest
antichain is $\{[n]\}$ and its event is the complete source key, whence
\begin{equation}
c_{\{[n]\}}(z;p)
=\log\frac{p(t\mid s_1,\ldots,s_n)}{p(t)}
=i(s_1,\ldots,s_n;t).
\label{eq:top}
\end{equation}
Define averaged atoms
\begin{equation}
\Pi_\beta^u(p)=\sum_{z:p(z)>0}p(z)\pi_\beta^u(z;p).
\label{eq:average}
\end{equation}
Linearity gives exact reconstruction:
\begin{equation}
\sum_\beta\Pi_\beta^{\mathrm{net}}(p)
=I(S_1,\ldots,S_n;T).
\label{eq:reconstruct}
\end{equation}

\paragraph{Finding F2 \evidence{P,R}.}
This is an algebraic theorem about distribution-dependent cumulatives.  It
does not require atoms to be independent target coordinates.

\subsection{What is actually nonnegative}

Makkeh--Gutknecht--Wibral prove, for each supported key,
\begin{equation}
\pi_\beta^+(z;p)\geq0,
\qquad
\pi_\beta^-(z;p)\geq0.
\label{eq:component-positive}
\end{equation}
This does not imply that $\pi_\beta^{\mathrm{net}}$ or its average is
nonnegative.

Independent rational checks found no violation of
Equation~\eqref{eq:component-positive}: all 12,869 binary two-source empirical
tables with $1\leq N\leq8$ (51,480 supported keys); all 4,844 binary
three-source tables with $1\leq N\leq4$ (15,504 keys); and forty seeded binary
four-source tables (1,060 keys, 166 atoms), with seed \texttt{1398296644}.
The first two searches were exhaustive within those bounds.  Ratios were
compared before logarithms.  Bounded checks support but do not replace the
proof.

\paragraph{Finding F3 \evidence{P,B}.}
No component-sign flaw was found.  Universal nonnegativity of \emph{net} atoms
would be false.

\section{What semantic choice is made?}

For propositions $A_1,\ldots,A_m$, a proposition $C$ follows from every $A_i$
if and only if it follows from their disjunction:
\begin{equation}
(\forall i,\;A_i\models C)
\quad\Longleftrightarrow\quad
(A_1\lor\cdots\lor A_m)\models C.
\label{eq:logic}
\end{equation}
The disjunction is, up to equivalence, the strongest proposition implied by
every branch.  Shared exclusions then identifies information shared by the
source statements with the likelihood-ratio information supplied by that
disjunction.  This is coherent, but it is a constitutive semantic postulate,
not a unique consequence of Shannon's axioms.  The parthood/logic isomorphism
establishes an indexing structure; it does not force every scientifically
useful redundancy measure to use Equation~\eqref{eq:cumulative}.

\paragraph{Finding F7 \evidence{P,R}.}
The logical claim is sound.  Uniqueness of the statistical meaning is not
established.

\subsection{There is no single fixed keyed auxiliary variable}

$A_\alpha(s)$ changes with the complete source key.  Consequently the average
in Equation~\eqref{eq:average} changes the event while weighting observations
by $p(s,t)$.  It is not generally $I(W;T)$ for one fixed random variable $W$.
The 2021 paper states this, rejects two naive channel readings, and supplies a
special metadata-masked channel interpretation.  An application claiming one
fixed ``shared signal'' needs an extra construction and equivalence proof.

\subsection{Independent COPY has no common deterministic signal}

Let $S_1,S_2$ be independent fair bits and $T=(S_1,S_2)$.  Bottom Sx
redundancy is
\begin{equation}
\Pi_{\{\{1\},\{2\}\}}^{\mathrm{net}}=\log\frac43>0.
\label{eq:copy-red}
\end{equation}
If $U=f_1(S_1)=f_2(S_2)$ almost surely, then $f_1(S_1)$ and $f_2(S_2)$ are
independent and equal.  Thus $U$ is independent of itself.  For every $u$,
\[
\Pr(U=u)=\Pr(U=u)^2,
\]
so every mass is zero or one and $U$ is constant.  Therefore $I(U;T)=0$ while
Equation~\eqref{eq:copy-red} is positive.

\paragraph{Finding F8 \evidence{X}.}
Sx redundancy is not G\'acs--K\"orner common information or the entropy of a
nonconstant statistic computable from each independent source.  Its commonality
is keyed and event-logical.

\section{Exact axiomatic trade-offs}

\subsection{Re-encoding invariance and target chain rule}

Coordinatewise bijections preserve equality events and all probabilities in
Equations~\eqref{eq:event}--\eqref{eq:components}.  They preserve every
cumulative and atom up to source-label permutation.  A joint invertible map
mixing sources changes the source partition and is not covered.

For $T=(T_1,T_2)$ and one fixed keyed event $A$,
\begin{equation}
\log\frac{p(t_1,t_2\mid A)}{p(t_1,t_2)}
=\log\frac{p(t_1\mid A)}{p(t_1)}
+\log\frac{p(t_2\mid t_1,A)}{p(t_2\mid t_1)}.
\label{eq:target-chain}
\end{equation}
The identity propagates through linear inversion and averaging, using the
conditional law in the second term.

\subsection{Exact identity and local-positivity failures}

For independent COPY, every key has $p(A)=3/4$, $p(t)=1/4$, and
$p(t,A)=1/4$.  The complete atom vector is
\begin{equation}
(R,U_1,U_2,\Syn)
=\left(\log\frac43,\log\frac32,\log\frac32,\log\frac43\right),
\label{eq:copy-atoms}
\end{equation}
and sums to $\log4$.  Identity redundancy would instead require
$R=I(S_1;S_2)=0$.

With the same sources but $T=S_2$, direct evaluation gives
\begin{equation}
(R,U_1,U_2,\Syn)
=\left(\log\frac43,-\log\frac43,\log\frac32,\log\frac43\right).
\label{eq:lp-failure}
\end{equation}
The negative unique atom has exact split
\begin{equation}
U_1^+=\log\frac32,
\qquad U_1^-=\log2,
\qquad U_1=-\log\frac43.
\label{eq:lp-components}
\end{equation}

\subsection{What arXiv:2512.16662v1 implies}

Matthias, Makkeh, Wibral, and Gutknecht prove that no general discrete
antichain PID satisfies all of net-atom local positivity (LP), target chain
rule (TCR), and coordinatewise invertible re-encoding invariance (REI).  They
also prove a related LP--identity--REI incompatibility.  SxPID has REI and TCR,
but not LP or identity.  Equations~\eqref{eq:copy-atoms} and
\eqref{eq:lp-failure} place it exactly on that trade-off.

This is a structural incompatibility, not an algebraic inconsistency of SxPID\@.
It does not imply that Sx fails reconstruction.  Version 1, dated 18 December
2025, is the exact version audited; arXiv still listed it as current when
rechecked on 25 July 2026.

\paragraph{Finding F6 \evidence{P,R,X,E}.}
Sx preserves TCR and REI by accepting signed net atoms.  A claim that it also
has universal net local positivity is false.

\subsection{XOR-source-copy executable witness}

Let $S_1,S_2$ be independent fair bits, $S_3=S_1\oplus S_2$, and
$T=(S_1,S_2,S_3)$.  The current kernel gives
\begin{align}
\Pi_{\{\{1\},\{2\}\}}
=\Pi_{\{\{1\},\{3\}\}}
=\Pi_{\{\{2\},\{3\}\}}&=\log\frac43,\\
\Pi_{\{\{1\},\{23\}\}}
=\Pi_{\{\{2\},\{13\}\}}
=\Pi_{\{\{3\},\{12\}\}}&=\log\frac98,\\
\Pi_{\{\{12\},\{13\},\{23\}\}}&=\log\frac{32}{27}.
\end{align}
All other atoms vanish, and
\[
\left(\frac43\right)^3\left(\frac98\right)^3\frac{32}{27}=4,
\]
so the atoms sum to $2\log2$.  This one unconditional witness need not display
a negative atom; the 2025 result concerns universal property combinations.

\section{Exact counterexample ledger}

\subsection{Negative redundancy and pairwise non-identifiability}

For independent fair $S_1,S_2$ and XOR target $T=S_1\oplus S_2$, every key has
$p(A)=3/4$, $p(t)=1/2$, and $p(t,A)=1/4$.  Thus
\begin{equation}
R_{\Sx}=\log\frac23<0.
\label{eq:xor-red}
\end{equation}
The vector
\[
\left(\log\frac23,\log\frac32,\log\frac32,\log\frac43\right)
\]
sums to $\log2$.  Compare this law with three independent fair bits
$S_1,S_2,T$.  Both have the same independent pairwise marginals $(S_1,T)$ and
$(S_2,T)$, but
\begin{equation}
R_{\Sx}^{\mathrm{XOR}}=\log\frac23,
\qquad R_{\Sx}^{\mathrm{independent}}=0.
\label{eq:pairwise}
\end{equation}
The full joint law, not pairwise marginals, identifies SxPID\@.

\subsection{An irrelevant source and source--target role asymmetry}

Add an independent fair $S_3$ to XOR without changing the target.  The bottom
event has $p(A)=7/8$, $p(t,A)=3/8$, hence
\begin{equation}
R_{\Sx}(S_1,S_2,S_3\to T)=\log\frac67,
\label{eq:irrelevant}
\end{equation}
which is larger than $\log(2/3)$.  Naive net source monotonicity is false.

For a separate exact role witness, assign probability $1/3$ to each row
\[
(S_1,S_2,T)\in\{(0,1,1),(1,0,0),(1,1,1)\}.
\]
Then
\begin{align}
R_{\Sx}(S_1,S_2\to T)&=\frac13\log\frac94,\label{eq:role-a}\\
R_{\Sx}(T,S_2\to S_1)&=\frac13\log\frac{27}{16}.
\label{eq:role-b}
\end{align}
Their difference is $\frac13\log(4/3)>0$.  PID is target-directed even though
ordinary mutual information is symmetric.

\subsection{No universal atomwise data processing}

Starting from independent COPY, coarsen the target deterministically to XOR
and then to a constant.  Bottom redundancy follows
\begin{equation}
\log\frac43\quad\longrightarrow\quad\log\frac23
\quad\longrightarrow\quad0.
\label{eq:coarse}
\end{equation}
The first step changes sign; the second removes all target information but
raises signed redundancy from a negative number to zero.  Coarsening one XOR
source to a constant also changes $\log(2/3)$ to zero.  Total mutual information
obeys data processing; individual signed Sx atoms do not inherit a universal
DPI\@.  Preprocessing changes the estimand.

\section{The 2026 structural-impossibility challenge}
\label{sec:lcr}

\subsection{Exact witness pair}

Lyu, Clark, and Raviv (arXiv:2604.03869v2, 14 April 2026) define two systems.
For the \emph{hat} system let
\[
x_1,x_2,x_4,x_5,x_7,x_8\stackrel{\mathrm{iid}}{\sim}
\operatorname{Bernoulli}(1/2),
\]
with
\[
x_3=x_1\oplus x_2,\qquad x_6=x_4\oplus x_5,
\qquad x_9=x_7\oplus x_8.
\]
For the \emph{tilde} system let
\[
x_1,x_2,x_4,x_5,x_7\stackrel{\mathrm{iid}}{\sim}
\operatorname{Bernoulli}(1/2),
\]
with
\[
x_3=x_1\oplus x_2,\qquad x_6=x_4\oplus x_5,
\qquad x_9=x_1\oplus x_5,
\qquad x_8=x_7\oplus x_1\oplus x_5.
\]
Both use
\begin{equation}
S_1=(x_1,x_4,x_7),\quad
S_2=(x_2,x_5,x_8),\quad
S_3=(x_3,x_6,x_9),\quad
T=(x_1,x_5,x_9).
\label{eq:lcr-system}
\end{equation}
The hat system has 64 equiprobable states and three independent target bits,
so $I(\widehat S;\widehat T)=3\log2$.  The tilde system has 32 states and
$x_9=x_1\oplus x_5$, so $I(\widetilde S;\widetilde T)=2\log2$.

\subsection{What Definition 6 and Lemma 4 assign}

Each target coordinate has the same minimal recovering antichain in both
systems:
\begin{equation}
\alpha(x_1)=\{\{1\},\{23\}\},\quad
\alpha(x_5)=\{\{2\},\{13\}\},\quad
\alpha(x_9)=\{\{3\},\{12\}\}.
\label{eq:access-labels}
\end{equation}
Their Lemma 4 stipulates
\begin{equation}
\Pi(\alpha)=H(U_\alpha).
\label{eq:descriptor}
\end{equation}
It assigns one bit to each of those labels and zero elsewhere in both systems.
In the tilde system the three coordinates are related.  Their stipulated
values sum to three bits while the target has two.  This is a recoverability
descriptor vector, not a whole-equals-sum-of-parts PID of $I(S;T)$ there.

The proof assigns $H(U_\alpha)$ directly; it does not invert an Sx cumulative,
and it does not prove independence among distinct $U_\alpha$.  Theorem 1 is
valid on its explicitly stated Definition 6 domain: the witness laws have the
same complete descriptor vector and unequal mutual information, so no single
map from that vector reconstructs mutual information for every law in that
domain.  Access labels and marginal component entropies omit the witness's
cross-coordinate relation.

\subsection{Exact SxPID result}

Independent exact event counting and the current Rust kernel instead give
eight nonzero net atoms.
\begin{center}
\small
\begin{tabularx}{\textwidth}{@{}Ycc@{}}
\toprule
\PidTableHeaderRow
Antichain family & Hat atom & Tilde atom \\
\midrule
$\{\{1\},\{2\},\{3\}\}$ & $\log(16/11)$ & $\log(8/5)$ \\
Each of $\{\{1\},\{2\}\}$, $\{\{1\},\{3\}\}$,
$\{\{2\},\{3\}\}$ & $\log(11/10)$ & $\log(15/14)$ \\
Each label in Equation~\eqref{eq:access-labels} & $\log(25/22)$ &
$\log(49/45)$ \\
$\{\{12\},\{13\},\{23\}\}$ & $\log(352/125)$ &
$\log(540/343)$ \\
\bottomrule
\end{tabularx}
\end{center}
The other ten atoms vanish.  Exact product identities establish reconstruction
before evaluating logarithms:
\begin{align}
\frac{16}{11}\left(\frac{11}{10}\right)^3
\left(\frac{25}{22}\right)^3\frac{352}{125}&=8,
\label{eq:hat-product}\\
\frac85\left(\frac{15}{14}\right)^3
\left(\frac{49}{45}\right)^3\frac{540}{343}&=4.
\label{eq:tilde-product}
\end{align}
Taking logs gives $3\log2$ and $2\log2$.  Eight of eighteen net atoms differ;
ten agree exactly.  The largest difference is
\begin{equation}
\log\frac{352/125}{540/343}
\approx0.58147874590342\ \text{nats}
=0.83889650309719\ \text{bits}.
\label{eq:max-delta}
\end{equation}
At each recovery label, the exact component atoms are
\[
\widehat\Pi^+=\log\frac{225}{176},\quad
\widehat\Pi^-=\log\frac98,
\qquad
\widetilde\Pi^+=\log\frac{49}{40},\quad
\widetilde\Pi^-=\log\frac98.
\]
The target relation changes informative increments although the access labels
stay fixed.

\subsection{What arXiv:2604.03869v2 does and does not imply}

\paragraph{Valid theorem on the stated Definition 6 domain \evidence{P,R}.}
The exhibited laws have equal complete vectors $(H(U_\alpha))_\alpha$ and
unequal mutual information.  Therefore no single function of that descriptor
vector reconstructs mutual information for every law in that domain.

\paragraph{Not established \evidence{X,E}.}
The paper does not establish that the two laws have equal SxPID atoms or that
SxPID fails MI reconstruction.  Equations~\eqref{eq:hat-product} and
\eqref{eq:tilde-product} show the opposite.

\subsubsection{Descriptor-factorization firewall}

The transfer condition does not require PID terminology.  Let $D(P)$ be a
descriptor vector, $A(P)$ an atom vector, and $J(P)$ the quantity to be
reconstructed.  Suppose the atom assignment factors through the descriptor:
\begin{equation}
A=g\circ D.
\label{eq:descriptor-factor}
\end{equation}
If two systems satisfy
\begin{equation}
D(P)=D(Q),\qquad J(P)\ne J(Q),
\label{eq:descriptor-collision}
\end{equation}
then $A(P)=A(Q)$, so no function $f$ can satisfy
$f(A(R))=J(R)$ for every system $R$.  Conversely,
\begin{equation}
D(P)=D(Q),\qquad A(P)\ne A(Q)
\quad\Longrightarrow\quad
\nexists g\; A=g\circ D.
\label{eq:factorization-refuted}
\end{equation}

This is the exact logical firewall.  The Lyu--Clark--Raviv witness proves a
reconstruction impossibility for atom assignments that accept
Equation~\eqref{eq:descriptor-factor}, using their recoverability-label and
component-entropy descriptor.  It does not transfer merely because another
assignment is indexed by the same antichain lattice.  A separate theorem must
establish Equation~\eqref{eq:descriptor-factor} for that assignment.

The SxPID calculation supplies the converse witness in
Equation~\eqref{eq:factorization-refuted}: the descriptors agree but eight Sx
atoms differ.  Therefore SxPID does not factor through that descriptor map.
The descriptor theorem remains valid; its operative premise is descriptor
sufficiency, not antichain indexing alone.

The generic firewall is kernel-checked in
\path{audit/formal/lean-foundational-sxpid/}\newline
\mbox{\path{PidDescriptorFactorization.lean}}.
Its three theorems have an empty Lean axiom inventory.  The pinned checker
\path{scripts/check-lean-descriptor-factorization.py} and fail-closed
three-mutation suite
\path{scripts/check-lean-descriptor-factorization-self-test.py} bind the result
to the recorded source, Lake manifest/configuration, toolchain identifier, and
reported Lean version.  The latter also kernel-checks three finite premise
countermodels.  Their deterministic outputs are
{\small\path{audit/evidence/foundational-sxpid-descriptor-factorization-lean-4.33.0.json}} and
{\small\path{audit/evidence/foundational-sxpid-descriptor-factorization-mutations-4.33.0.json}}.
These are reproducibility checks, not an authenticity or binary-attestation
claim about the installed Lean/mathlib toolchain.

Phrases such as ``regardless of any axioms'' and ``any multivariate
information decomposition that relies solely on the antichain lattice'' are
therefore safe only with the descriptor-factorization premise explicit.
Distribution-dependent signed M\"obius atoms, including SxPID, do not accept
that premise in this witness.

\paragraph{Preserved insight.}
Access labels alone do not encode every relation among semantically named
target components.  Relation-aware metadata may be useful, but it cannot
replace Sx atoms with nonadditive coordinate entropies and remain a
decomposition of mutual information.

\paragraph{Finding F11 \evidence{P,R,X,E}.}
The witness is a genuine warning against access-label/marginal-entropy
descriptors as a complete representation.  It is not a contradiction of
shared-exclusions PID\@.

\subsection{Three independent executable and formal routes}

The standard-library checker
\path{audit/tools/foundational_sxpid/check_lcr_relation_}\newline
\mbox{\path{witness.py}}
does not import \texttt{pid-rs}.  It:
\begin{enumerate}
  \item verifies the generated rows match the paper's fair-bit/XOR
        constructions and source/target index sets;
  \item checks Definition~6 conditions (i), (ii), and (iii) exactly on both
        finite laws, including
        mutual rather than pairwise within-source independence and all
        twenty-four target-component/source-group cases per system;
  \item derives all minimal recovering antichains and both complete Lemma 4
        descriptor vectors from the rows, proving those vectors equal;
  \item generates all eighteen nonempty antichains of the seven nonempty
        three-source subsets;
  \item independently implements Equation~\eqref{eq:order};
  \item counts all keyed equality-event unions and target intersections;
  \item performs exact M\"obius inversion in formal prime-exponent logarithms;
  \item verifies all local component atoms are constant across these two
        equiprobable supports;
  \item derives the eight differing atoms, ten equal atoms, and products 8 and
        4; and
  \item kills union-to-intersection, row-derived Lemma 4 descriptor substitution,
        and cross-relation mutations.
\end{enumerate}
Its full output is
\path{audit/evidence/foundational-sxpid-lcr-exact-audit.json}, bound to the
checker, generated rows, Rust regression, and Rust kernel by SHA-256 digests.
The independent Rust route is
\path{crates/pid-core/tests/sxpid_relation_witness.rs}; it calls the public
categorical kernel, checks every displayed atom and component, enforces the
integer product identities, and rejects descriptor substitution.

The third route is the generic Lean factorization firewall above.  Its proof
has no axioms in the kernel inventory.  The self-test kills removal of the
factorization premise, replacement of the unequal-quantity premise by
equality, and replacement of the unequal-atom premise by equality.  Three
finite Bool/Unit countermodels show why the altered premises cannot support the
original conclusions: equal descriptors alone need not equalize atoms, equal
quantities can admit reconstruction, and equal atoms can admit factorization.
This route proves the transfer logic, not the concrete Sx event calculation;
the exact-rational and Rust routes supply that separate binding.

At generation, the bound SHA-256 values were:
\begin{itemize}
\small
  \item checker:
        \path{8244bd6a3a3a590f8a8da9f63ce5d5114d3453cb4623bde70ce02c7ddbeec96a};
  \item Rust regression:
        \path{e11076dfd0bc8e8b3f3565128ed87c2cabe03c0f25b966e76fe925a09e3451ed};
  \item Rust kernel:
        \path{f08b56bac473474b39dd2cf2f09c5d13ccda025825a011ee1870f6f2979ff98a}.
\end{itemize}

\section{Support, continuity, and estimation}

Every probability in Equation~\eqref{eq:cumulative} is positive at a supported
key.  If that key disappears, its local quantity has no operational occurrence;
along rare-key sequences a local term can diverge like $-\log p(z)$.  The 2021
differentiability theorem is an interior theorem on a fixed finite simplex, not
pointwise continuity at disappearing keys.

A separate project theorem proves that support-restricted \emph{averaged}
cumulatives and full-lattice atoms on one fixed finite alphabet extend
continuously through support changes.  A vanishing weight can control the local
log divergence:
\[
\text{local value may diverge}
\quad\text{while}\quad p(z)\pi_\alpha(z;p)\longrightarrow0.
\]
This does not make absent-key pointwise quantities continuous.

For a known finite alphabet, the full joint population law identifies SxPID\@.
Plug-in consistency does not supply support discovery, useful sparse-law rates,
protection against adaptive quantization, generic post-selection/dependence
inference, binary64 sign certificates, or causal identification.  Reports need
the estimand and source partition, alphabet and occupancy, sampling contract,
preprocessing split, both components, sign status, reconstruction residual,
sensitivity analysis, and an explicit noncausal boundary.

\paragraph{Finding F4/F10 \evidence{R}.}
Exact-real continuity and fixed-alphabet consistency are valid stability
results; they do not establish practical calibration in every data regime.

\section{Interpretive and causal limits}

A net atom is a M\"obius increment in a likelihood ratio.  Its sign compares
informative and misinformative increments.  Without additional assumptions it
does not mean beneficial mechanism, causal redundancy/synergy, correct sensor,
deception, intent, replaceability, intervention robustness, safety margin, or
decision correctness.

SxPID is a functional of an observational joint law.  Distinct causal graphs
can induce the same law and SxPID\@.  No statistic of that law can distinguish
them absent causal assumptions or interventions.  The published operational
interpretation also assumes a particular metadata-masked event channel.  A
deployment needs a loss model, calibration, sequential error policy, drift and
out-of-distribution rules, and abstention.  SxPID can be evidence, not an
authorization rule.

\paragraph{Finding F12 \evidence{R}.}
Causal or authority-grade conclusions require a separate theorem connecting
the observational Sx estimand to the needed intervention or decision quantity.

\section{Validity-regime contract}

Categorical SxPID is defensible when variables are categorical (or quantization
is explicitly the estimand), the scientific source partition is fixed, the
question is shared logical-exclusion information, signed atoms are acceptable,
both components are retained, the full joint law is estimable under a declared
sampling regime, dependence/drift/selection are separately controlled,
cancellation and reconstruction are checked, and interpretation remains
observational absent a causal bridge.

Another object is required for a common deterministic random variable, only
nonnegative parts, identity redundancy, arbitrary local-processing
monotonicity, source--target symmetry, pairwise identification, causal effect,
unsupported high-dimensional continuous PID, or an authority-grade probability
of correctness.  Those are different mathematical requirements.

\section{Formal-verification consequences}

\subsection{Exact specification and executable refinement}

Formalize finite alphabets; supported keys; equality events; collections and
antichains; Equations~\eqref{eq:event}--\eqref{eq:components}; the order and
containment; concrete M\"obius inversion; the top-node identity; averaged
reconstruction; component positivity; coordinatewise REI; and target chain.
Formalize identity, LP, pairwise-identification, role-symmetry, irrelevant-source,
and net-DPI counterexamples as theorems.  These negative specifications prevent
silent future claim strengthening.

Prove that canonical row counting equals the empirical law; each event count
equals its predicate; the Rust order matches Equation~\eqref{eq:order}; every
cumulative reconstructs; specialized/general paths agree; admitted integer
count products cannot overflow; and reports preserve signed components.

\subsection{Certified numerical and semantic layers}

Counts give rational pre-log quantities.  Directed-rounding intervals can
certify every log, atom, sum, and reconstruction.  Report certified positive,
certified negative, or unresolved sign rather than clamping.

Keep semantic boundaries explicit: disjunction as shared logical consequence
is an accepted postulate; uniqueness among all redundancy meanings is not a
theorem; atoms are not common deterministic coordinates or nonnegative set
sizes; observational atoms are not causal effects.  Formal methods can prove a
precise model correct, but cannot make an unchosen semantics uniquely true.

\section{Correctness-preserving research directions}

\begin{enumerate}
  \item \textbf{Relation annotations, not replacement atoms \evidence{O}.}
        Attach target multi-information, algebraic constraints, or declared
        factor graphs as metadata while preserving
        Equation~\eqref{eq:reconstruct}.  Do not add them again as atoms without
        a no-double-counting theorem.
  \item \textbf{Semantic uniqueness class \evidence{O}.}  Find the weakest
        logical, epistemic, decision-theoretic, or Shannon axioms under which
        disjunction is uniquely selected.
  \item \textbf{Coarse-graining calculus \evidence{O}.}  Classify maps that
        preserve the keyed event sigma-algebra, merge atom families in a known
        way, or preserve signs.  Keep Equations~\eqref{eq:irrelevant} and
        \eqref{eq:coarse} as negative controls.
  \item \textbf{Operational channel validation \evidence{O}.}  Test whether
        real observers receive the assumed metadata-masked channel.  Leakage or
        learning across uses changes the operational estimand.
  \item \textbf{Finite-sample signed certification \evidence{O}.}  Combine
        valid concentration, closed-simplex continuity, and intervals while
        separating model, sampling, selection, and numerical error.
\end{enumerate}

\section{Final claim ledger}

\small
\begin{longtable}{@{}P{0.44\textwidth}P{0.20\textwidth}P{0.25\textwidth}@{}}
\toprule
\PidTableHeaderRow
Claim & Decision & Evidence \\
\midrule
\endfirsthead
\toprule
\PidTableHeaderRow
Claim & Decision & Evidence \\
\midrule
\endhead
Keyed cumulative is a valid local likelihood ratio & Established &
\evidence{P,R}, Equations~\eqref{eq:event}--\eqref{eq:components} \\
Atoms reconstruct mutual information & Established & \evidence{P,R,E},
Equations~\eqref{eq:zeta}--\eqref{eq:reconstruct} \\
Component atoms are nonnegative & Established/\allowbreak bounded-tested & \evidence{P,B} \\
Net atoms are nonnegative & False & \evidence{X,E},
Equations~\eqref{eq:lp-failure}, \eqref{eq:xor-red} \\
Coordinatewise REI and target chain hold & Established & \evidence{P,R} \\
Identity holds & False & \evidence{X,E}, Equation~\eqref{eq:copy-atoms} \\
Sx redundancy is common deterministic information & False & \evidence{X} \\
Pairwise marginals identify Sx redundancy & False & \evidence{X},
Equation~\eqref{eq:pairwise} \\
Net atoms obey arbitrary processing DPI & False & \evidence{X},
Equation~\eqref{eq:coarse} \\
Source/target roles are interchangeable & False & \evidence{X},
Equations~\eqref{eq:role-a}--\eqref{eq:role-b} \\
Disjunction semantics is logically coherent & Established & \evidence{P,R} \\
Disjunction is uniquely forced by Shannon theory & Not established &
Constitutive-semantic gap \\
The Lyu systems have identical SxPID atoms & False & \evidence{X,E},
Equations~\eqref{eq:hat-product}--\eqref{eq:tilde-product} \\
Lyu descriptors expose lost target relations & Established in their model &
\evidence{P,R} \\
Descriptor-collision impossibility transfers to every antichain PID & False
without a factorization theorem & \evidence{P,R,X,E},
Equations~\eqref{eq:descriptor-factor}--\eqref{eq:factorization-refuted} \\
Lyu refutes all lattice PID including Sx & Not established for Sx &
Descriptor factorization fails and three routes agree \\
SxPID is causal/authority-grade by itself & False/unsupported & Observational
identification boundary \\
Categorical SxPID is a defensible signed event-logical estimand & Yes under the
regime contract & Entire audit \\
\bottomrule
\end{longtable}
\normalsize

\section{Reproducibility record}

The Rust witness was evaluated against commit
\path{6e29b3f5badc337a58eb821babc41641037e15d0}.  The audited kernel SHA-256 was
\begin{center}
\texttt{f08b56bac473474b39dd2cf2f09c5d13ccda025825a011ee1870f6f2979ff98a}.
\end{center}
The environment was Rust/Cargo 1.96.0 and Python 3.14.6.  The permanent commands
are
\begin{verbatim}
python3 -I -S audit/tools/foundational_sxpid/check_lcr_relation_witness.py \
  --write-evidence audit/evidence/foundational-sxpid-lcr-exact-audit.json
cargo test --locked -p pid-core --test sxpid_relation_witness
python3 -I -S scripts/check-lean-descriptor-factorization.py
python3 -I -S scripts/check-lean-descriptor-factorization-self-test.py
\end{verbatim}
\smallskip
\noindent At creation both Rust tests, the independent exact checker, and the Lean
factorization checks passed.  The
rational/formal-log route derives event counts, order, atoms, and mutations
without invoking the Rust implementation.  The Rust route independently
exercises the public event/lattice kernel.  The Lean route proves the abstract
transfer firewall while explicitly leaving the concrete descriptor/Sx binding
to the two executable routes.

\phantomsection
\section*{Primary sources}
\addcontentsline{toc}{section}{Primary sources}
\begin{enumerate}
  \item A. Makkeh, A. J. Gutknecht, and M. Wibral,
  ``Introducing a differentiable measure of pointwise shared information,''
  \emph{Physical Review E} 103, 032149 (2021),
  \href{https://doi.org/10.1103/PhysRevE.103.032149}{\nolinkurl{doi:10.1103/PhysRevE.103.032149}},
  \href{https://arxiv.org/abs/2002.03356}{\nolinkurl{arXiv:2002.03356v5}}.
  \item A. J. Gutknecht, M. Wibral, and A. Makkeh,
  ``Bits and Pieces: Understanding Information Decomposition from Part-whole
  Relationships and Formal Logic,'' \emph{Proceedings of the Royal Society A}
  477, 20210110 (2021),
  \href{https://doi.org/10.1098/rspa.2021.0110}{doi:10.1098/rspa.2021.0110},
  \href{https://arxiv.org/abs/2008.09535}{arXiv:2008.09535}.
  \item P. H. Matthias, A. Makkeh, M. Wibral, and A. J. Gutknecht,
  ``Novel Inconsistency Results for Partial Information Decomposition,''
  \href{https://arxiv.org/abs/2512.16662}{arXiv:2512.16662v1},
  18 December 2025.  Version 1 was audited and remained current on arXiv when
  rechecked on 25 July 2026.
  \item A. Lyu, A. Clark, and N. Raviv,
  ``Structural Impossibility of Antichain-Lattice Partial Information
  Decomposition,''
  \href{https://arxiv.org/abs/2604.03869}{arXiv:2604.03869v2},
  14 April 2026.
  \item P. L. Williams and R. D. Beer, ``Nonnegative Decomposition of
  Multivariate Information,''
  \href{https://arxiv.org/abs/1004.2515}{arXiv:1004.2515} (2010).
\end{enumerate}

\end{document}
